% !TEX root = ../../hw03.tex
% Homework 03, question 7.
\begin{proof}[Solution]
  Write $x=(x-y)+y$. The triangle inequality gives
  \[
    |x|\leq|x-y|+|y|,
  \]
  and therefore $|x|-|y|\leq|x-y|$. Interchanging $x$ and $y$ gives
  $|y|-|x|\leq|y-x|=|x-y|$. Combining the two inequalities, we obtain
  \[
    -|x-y|\leq|x|-|y|\leq|x-y|.
  \]
  This is exactly
  \[
    \bigl||x|-|y|\bigr|\leq|x-y|.
  \]

  For equality, first suppose $xy\geq0$. If both numbers are
  nonnegative, then
  \[
    \bigl||x|-|y|\bigr|=|x-y|.
  \]
  If both are nonpositive, then $|x|=-x$ and $|y|=-y$, so
  \[
    \bigl||x|-|y|\bigr|=|-x+y|=|x-y|.
  \]
  These cases also include either number being zero.

  If $xy<0$, the numbers have opposite signs and neither is zero.
  Then $|x-y|=|x|+|y|$. Since both $|x|$ and $|y|$ are positive,
  \[
    \bigl||x|-|y|\bigr|<|x|+|y|=|x-y|.
  \]
  Thus equality holds exactly when $xy\geq0$. Taking absolute values
  can be pictured as folding the number line at zero. Points on the
  same side keep their separation, while points on opposite sides
  move closer together. \autoref{fig:hw03:reverse-triangle-distance}
  illustrates this folding.
  \begin{marginfigure}
    \centering
    % Original folding analogy for the reverse triangle inequality.
% Background on absolute value as distance from zero and geometric inequalities:
% https://danaernst.com/IntroToProofViaIBL/sec_AxiomsRealNumbers.html
% Both rows use the same scale; the illustrated case is y < 0 < x < |y|.
\begin{tikzpicture}[interval diagram,x=7.5mm,y=1cm]
  \draw[<->,gray!60] (-2.45,0) -- (2.45,0);
  \foreach \position/\name in {-2/y,0/0,1/x} {
    \node[interval closed] at (\position,0) {};
    \node[interval label,below=5pt] at (\position,0) {$\name$};
  }
  \draw[<->] (-2,0.65) -- (1,0.65)
    node[midway,above=5pt,interval label] {$|x-y|$};

  % Move the negative point onto the positive half-line without crossing
  % either distance arrow. Keep the annotation in the open part of the bend.
  \draw[->,gray] (-2,-0.55)
    .. controls (-2,-2.2) and (2,-1.2) .. (2,-3.0);
  \node[interval label,anchor=west] at (-2.35,-2.25) {fold at $0$};

  \draw[->,gray!60] (0,-3.5) -- (2.45,-3.5);
  \foreach \position/\name in {0/0,1/{|x|},2/{|y|}} {
    \node[interval closed] at (\position,-3.5) {};
    \node[interval label,below=5pt] at (\position,-3.5) {$\name$};
  }
  \draw[<->] (1,-4.35) -- (2,-4.35)
    node[midway,below=5pt,interval label] {$\bigl||x|-|y|\bigr|$};
\end{tikzpicture}

    \caption[Folding the number line at zero.]{Schematic folding by $t\mapsto|t|$.
      Both rows use the same scale. Here $y<0<x<|y|$, and the gap shrinks
      from $|x-y|$ to $\bigl||x|-|y|\bigr|$. Points on the same half-line
      keep their separation.}
    \label{fig:hw03:reverse-triangle-distance}
  \end{marginfigure}
\end{proof}
